\subsection{Oppfølging og kontroll}
\label{subsec:oppfolging}

Aktivitetsprogrammet strekker seg over en pilotfase og en 24~måneders
langtidsplan som til sammen krever aktiv styring, ikke bare planlegging.
\textcite[697]{kotlerkeller2016} definerer markedskontroll som prosessen
der selskaper vurderer effekten av egne markedsaktiviteter og gjør
nødvendige justeringer. Poenget er at kontroll uten korrigerende tiltak
har liten verdi, og at kontroll uten tydelig eierskap ender som en
rapporteringsrutine ved siden av driften. Seksjonen beskriver først
kontrollen i pilotfasen, deretter kontrollen i langtidsplanen, og til
slutt ansvarsfordelingen i en samlet oversikt.

\subsubsection{Kontroll i pilotfasen}

Pilotprosjektet bør følge de tre KPI-typene som ble introdusert i
pilotbeskrivelsen, nemlig operative indikatorer (leveringstid, punktlighet,
avviksrate, vær- og sikkerhetsrelaterte kanselleringer), markedsmessige
indikatorer (førstegangskjøp, konvertering, prisaksept, gjenkjøp,
CAC/LTV) og aksept- og trygghetsindikatorer (opplevd trygghet og
enkelhet, naboaksept, bekymringer rundt støy og personvern). Disse måles
kontinuerlig gjennom hele pilotperioden, ikke bare ved avslutning, og
eierskapet er fordelt som vist i ansvarsfordelingen i
tabell~\ref{tab:ansvar-kontroll}.

For at KPI-ene skal styre atferd og ikke bare dokumentere den, bør hver
indikator ha en forhåndsdefinert terskelverdi og et definert
handlingsrom dersom terskelen ikke nås. Hvis opplevd trygghet etter
tillitsfasen ligger vesentlig under forventet nivå, bør ikke Zipline
uten videre gå inn i verdifasen, men forlenge tillitsarbeidet, justere
kommunikasjonen rundt godkjenning og sikkerhet, og håndtere konkrete
naboinnvendinger før neste fase starter. Hvis punktligheten svikter
jevnt på bestemte leveringspunkter, er det ikke markedsføringen som skal
justeres, men ruteplanleggingen eller basedriften. Hvis gjenkjøpsraten i
vanedannelsesfasen er lav til tross for introduksjonstilbud og
førstegangsrabatter, er det verdiforslaget eller appopplevelsen som må
undersøkes, ikke kommunikasjonen. Avvik betyr ikke automatisk at piloten
har mislyktes, men utløser en diagnose av årsaken og et konkret tiltak
før arbeidet går videre.

Overgangene mellom pilotens tre faser (tillit, verdi og vanedannelse)
bør dermed være eksplisitte beslutningspunkter, ikke automatiske
overganger på kalenderen. Pilotprosjektleder har det operasjonelle
eierskapet til disse overgangene og kan velge å forlenge en fase,
revidere tiltakene eller gå videre som planlagt. Beslutningen om å
avslutte piloten og gå til en bred Bærum-lansering er derimot et
strategisk beslutningspunkt som eies av ledelsen, fordi den utløser
betydelige investeringer og en ny risikoeksponering.

\subsubsection{Kontroll i den langsiktige planen}

Over den 24~måneder lange ekspansjonsplanen fungerer tre kontrollnivåer
med ulik frekvens sammen. Hvert nivå svarer på en distinkt type spørsmål
og er knyttet til et distinkt eierskap, slik at ansvaret for
korrigerende tiltak havner riktig sted.

\paragraph{Månedlig driftskontroll.}
På månedsbasis måles operative og markedsmessige nøkkeltall per aktivt
marked: leveringstid, punktlighet, avviks- og kanselleringsrate,
kundetilfredshet, CAC, konvertering og gjenkjøp. Operasjonell ansvarlig
eier driftsrelaterte indikatorer; markedsansvarlig eier kampanje- og
konverteringsmetrikkene. Formålet er å fange operative problemer før de
rammer kundeopplevelsen bredt, og å oppdage kampanjer som ikke virker
før budsjettet brennes. Avvik bør diagnostiseres etter kilde slik at
tiltaket treffer problemet: svikter leveringstiden i Bærum konsekvent i
rushtiden, er det baselokasjonen eller ruteplanleggingen som må
revideres, ikke kampanjebudskapet. Faller gjenkjøpsraten etter
introduksjonskampanjen, er det verdiforslaget eller appopplevelsen som
må undersøkes, ikke leveringstiden.

\paragraph{Kvartalsvis lønnsomhetskontroll.}
Hvert kvartal brytes resultatene ned på segment, område og kundegruppe
for å identifisere hvor Zipline tjener og taper penger. Dette brukes til
å reallokere ressurser mellom fasene. Dersom gjenkjøpsraten er vesentlig
høyere i Bærum enn i Vestre Aker, bør markedsbudsjettet vektes tyngre
mot Bærum selv om den geografiske planen i utgangspunktet forutsetter
parallell drift. Dersom én kundegruppe innenfor et segment viser markant
høyere livstidsverdi enn de øvrige, bør målrettingen strammes inn mot
denne gruppen. Eierskapet ligger hos økonomiansvarlig, mens beslutningen
om omfordeling av markedsbudsjett tas i dialog med markedsansvarlig. Den
kvartalsvise analysen leverer også datagrunnlaget som fase~3 allerede
forutsetter (spørsmålet om tilgjengelighet eller konkurransesituasjon er
den dominerende lønnsomhetsvariabelen), uten at selskapet må vente til
måned~12 for å se mønsteret. Konkrete tall og følsomheter tallfestes
i den økonomiske analysen.

\paragraph{Strategisk revisjon ved faseovergangene.}
Ved måned~6, 12, 18 og 24 gjennomføres en bredere strategisk
gjennomgang som ikke bare vurderer om planen går som tenkt, men om
planen fortsatt er den riktige. Gjennomgangen tester tre forhold: om
segmenteringen fortsatt treffer, eller om data peker mot andre
geografier eller kundegrupper; om posisjoneringen som raskere og
rimeligere faktisk dokumenteres i opplevd kundeverdi, eller om
konkurransedynamikken krever en justert differensiering; og om de tre
antagelsene planen bygger på (ingen Wing/Manna-inntreden, ingen raskere
åpning av urban dronedrift, ingen vesentlige vinterdrevne avbrudd)
fortsatt holder. Hvis en antagelse faller, er det ikke en månedlig
justering som trengs, men en revisjon av planens overordnede logikk. Det
regulatoriske og politiske arbeidet inngår også i den strategiske
revisjonen, ettersom endringer her direkte påvirker hvilke faser som er
gjennomførbare og i hvilken rekkefølge. Eierskapet ligger hos ledelsen,
og hver faseovergang er et bevisst beslutningspunkt, ikke en automatisk
kalendermessig overgang.

\subsubsection{Ansvarsfordeling}

Tabell~\ref{tab:ansvar-kontroll} samler ansvarsfordelingen for
kontrollpunktene på tvers av aktivitetsprogrammet. Rolletitlene er
indikative og kan tilpasses den interne organiseringen Zipline etablerer
i det norske markedet, men prinsippet er at hvert kontrollpunkt bør
ligge nærmest den som faktisk har tilgang til å justere det som måles.
På den måten blir kontrollen et reelt styringsverktøy for den toårige
planen, og ikke en rapporteringsrutine ved siden av driften.

\begin{table}[H]
\centering
\caption{Ansvarsfordeling for kontrollpunktene i aktivitetsprogrammet}
\label{tab:ansvar-kontroll}
\renewcommand{\arraystretch}{1.3}
\begin{tabularx}{\textwidth}{@{} >{\bfseries\raggedright\arraybackslash}p{5.6cm} >{\raggedright\arraybackslash}X @{}}
\toprule
\rowcolor{tableheader}
\color{white}\textbf{Kontrollpunkt} &
\color{white}\textbf{Eierskap og beslutningsrom} \\
\midrule
Operative KPI-er (leveringstid, punktlighet, avviksrate) &
Operasjonell ansvarlig. Kan justere ruteplanlegging, basedrift og
flåtedisponering. \\
Markedsmessige KPI-er (CAC, konvertering, gjenkjøp) &
Markedsansvarlig. Kan justere kampanjer, kanalmiks og kreativ retning
innenfor gitt budsjett. \\
Aksept- og trygghetsindikatorer &
Kommunikasjonsansvarlig i dialog med pilotprosjektleder. Kan justere
sikkerhetskommunikasjon, naboarbeid og lokal PR. \\
Kvartalsvis segmentlønnsomhet &
Økonomiansvarlig, i dialog med markedsansvarlig om omfordeling av
markedsbudsjett mellom fasene. \\
Faseovergangene i piloten (tillit $\to$ verdi $\to$ vane) &
Pilotprosjektleder. Go/no-go til neste fase, med mulighet til å forlenge
eller revidere fasen. \\
Faseovergangene i langtidsplanen (måned 6, 12, 18, 24) &
Ledelsen. Strategisk go/no-go og eventuell revisjon av planens
overordnede logikk. \\
Regulatorisk arbeid (Luftfartstilsynet, EASA) &
Dedikert regulatorisk ansvarlig gjennom hele perioden, som beskrevet i
langtidsplanen. \\
Politisk arbeid og offentlig finansiering &
Ledelsen i samråd med regulatorisk ansvarlig. \\
Partnerrelasjoner (Coop Obs, Jordbærpikene, restauranter i Bærum og
Holmenkollen-området) &
Partneransvarlig. Eier både rammevilkår og daglig koordinering. \\
Ambassadørprogram i fase 2 (Holmenkollen, Slemdal, Vinderen) &
Dedikert relasjonsansvarlig med lokal forankring. \\
\bottomrule
\end{tabularx}
\end{table}
